\section{Uniform parameters and the actual Type-I height}
\label{sec:numerical-parameters}

\status{NP1--NP9 supply the literal standard, segment, rounded-weight and
projection numerical records, including the shifted comparison with $t_I$.
These finite inequalities do not assert the existence of the global
arithmetic data to which the local theorem is to be applied.}

\paragraph{NP1--NP2: a uniform standard package.}
Write $k=m+1$, $N=|TT|$, $Q=p^{G'}$, $\lambda=C_0N_1^2\log a$,
$\epsilon=N_1\lambda/k$, and
\[
 E_0=C_{E1}\left(\frac a{p^2}+\frac {b'}p+
              \frac{\lambda N_1a}{pk}+\frac{\lambda b'}k\right)
       +C_{E2}k^2N_1Q^2c.
\]
For $H\ge1$, $C_c\ge0$, and $C_0,C_{E1},C_{E2}\ge1$, set
\[
 T_2=1000(C_{E1}+C_{E2}+1)(C_c+1)H.                 \tag{NP.1}
\]
The literal $\mathrm{D1num}$, extended-window hypothesis $\mathrm{HIext}$,
and $\mathrm{CBound}(C_c)$ imply
\[
                  \mathrm{Std}(C_0,C_\delta,C_{E1},C_{E2},H,X)
\]
whenever
each of $C_{Yb},C_\delta,C_k,C_D,C_b$ is at least $T_2$.
The report lists every field of these records. The substantive estimates
are as follows.

Put $F=c_*+Nc_{\rm cl}+n'c_{\rm cl}$. The depth hypothesis gives
$Nk\log_pa\le D_{\rm LP}$, whence $c\le2C_cF$.
The size hypothesis and $C_b\ge2C_cH,H$ give
\[
 Hk^3NQ^4cY\le b',\qquad Hk^3NQ^4Y\le b'.        \tag{NP.2}
\]
Set $A=N_1\lambda/k^2$. Each of
\[
 a/p^2,\quad b'/p,\quad\lambda N_1a/(pk),\quad
 \lambda b'/k,\quad k^2N_1Q^2c,\quad Q
\]
is at most $aA$: use $p\ge k^2$, $kb'\le a$,
$k^3Q^2c\le b'$, $k^3Q\le b'$, and $N_1,\lambda\ge1$.
Since $k\ge4C_{E1}+C_{E2}+1$, this proves
\[
                         (E_0+Q)/a\le\epsilon.    \tag{NP.3}
\]
If $\eta=(b_h-b_l)/b'$, then
$\delta=C_\delta\epsilon+2\eta$ and
$\beta=\eta+2\lambda/k$. Thus
$40\epsilon\le\delta<1/6$ and
$0\le\beta\le\delta/2-18\epsilon$.
Finally $2J_0=2\lambda b'/k\le a/(kY)$ enlarges the original windows
to exactly the standard windows. All other fields follow directly or
by removing factors at least one.
The quantifier order is preserved: choose $C_c$, then $C_0\ge1$,
then $C_{E1},C_{E2},C_m\ge1$, and only then the five constants
bounded below by (NP.1). In particular $H=4$, needed by RD, is uniform.

\paragraph{NP3: the actual segment satisfies the size conditions.}
For a segment put $r=m$, $s\le N_1$, and $P_i=p^{G-d_i}$.
The extended hypotheses give the source record $\mathrm{SegFacts}$:
\[
 \begin{gathered}
 p\ge(r+1)^2Y^2,\quad 2s+2\le r,\quad
 p\le P_i\le Q_j\le Q,\quad (r+1)P_i\le b',\\
 4\lambda<r+1,\quad (r+1)b'\le a/4,\quad
 b_T\le\tfrac52b',\quad 0\le x_a\le2J_0.
 \end{gathered}                                      \tag{NP.4}
\]
Here $C_{Yb}\ge2,C_b\ge1,C_D,C_{E1},C_{E2}\ge0$ suffice.
Indeed $p\ge k^2\ge4k$ gives $kb'\le a/4$, while the size
inequality gives $kQ\le b'$ and $2Y_bN\lambda<k$ gives $4\lambda<k$.
For the multiple-tie model of TP, a nonterminal letter degree $d_l$
satisfies $d_lP_i\le2a$ and the terminal degree is at most $2a$.
Consequently all letter and tie degrees are at most $2a$ and the
sum of letter degrees is at most $4a\le8a$.
The two bounds
\[
 S_iP_i\le\tfrac52b'+(2s+1)2a/p,\qquad
 n_{0i}P_i\ge a-kP_i-\tfrac52b'
\]
give $S_i\le n_{0i}$ because $p\ge4s+8$ and $pb'<a$.
They apply when all tie coordinates are simultaneously at their
compatible endpoints, by SW6.
Under $C_0\ge1$ and the NP2 thresholds, the remaining bounds in
SW5--SW7 follow from $kY^2b'\le a/2$, $p\ge k^2Y^2$,
$a/(kY)\le|W|\le a+1$, and SW12.
Thus TP11 applies to the actual model with $C_M=1$, including its
logarithmic degree hypotheses; SW5--SW9 need no additional numerical
assumption. The coarse rank input is supplied separately by RU.

\paragraph{NP4: rounding the radius without changing the hypotheses.}
For the final Type-I statement distinguish the real radius $a_r$ from
$a=\lceil a_r\rceil$, and put $L=\log_pa_r$.
For a fixed positive constant tuple $K$, define
\[
 \begin{split}
 T_{\rm LP}&=\max\{108,\,8C_A+93,\,1,\,C_0^{-2}\},\\
 T_M&=\max\{T_{\rm LP},\,4K.CH+1,\,2C_m+24\}.
 \end{split}                                         \tag{NP.5}
\]
Assume the literal $\mathrm{HypI}(C_0,C_H)$, with $C_H\ge T_M$.
Its elementary consequences are $k\ge5,p\ge25,b'\ge8,a_r\ge200$,
$N,N_1^2\le k$, and $L,\log a_r,\lambda(a_r)\ge1$.
Since $a_r\le a<a_r+1$,
\[
 0\le\log a-\log a_r\le1/200,\quad
 \lambda(a)\le2\lambda(a_r),\quad
 \lambda(a)\log a\le2\lambda(a_r)\log a_r.          \tag{NP.6}
\]
Replace $a_r$ by $a$, $Y_b$ by $Y_b/2$, and the constant $C_H$ by
$K.CH$. The resulting data still satisfy $\mathrm{HypI}$.
For the delicate clause, write $\delta_r,\delta_a$ for the
old and new excluded-band widths. Then
$\delta_a+1/k\le\delta_r$, while $k\le a_r$ and $\delta_a<1/6$.
These imply
\[
 [a(1-2\delta_a),a(1+2\delta_a)]
 \subseteq[a_r(1-2\delta_r),a_r(1+2\delta_r)].        \tag{NP.7}
\]
The lower endpoint uses monotonicity. For the upper endpoint,
the rounding error is at most $1+2\delta_a<4/3$, whereas the
additional allowance is $2a_r/k\ge2$.
The other clauses use (NP.6), with factors two or four accounted for
by (NP.5). In particular the forbidden odd powers remain forbidden.

\paragraph{NP5--NP6: depths, windows and multiplicity.}
Put $\ell=1+\log_p(8a+m)$, $u_0=\lceil(k+1)\ell\rceil$,
$\kappa_1=\lceil(k-1)^2\ell\rceil$, $n_s=\mathrm{numSeg}(TT)$,
and $v=\log_pa$. The source depths are
\[
 \begin{split}
 D_0&=n_s(C_AY^2N_1^2k^2\ell+4Nkv)+8Nkv,\\
 D_1&=D_0+4Nu_0.
 \end{split}
\]
Here $\ell\le3L$, $v\le2L$. With $P=NY^2k^3L$, the four contributions
to $D_1$ are bounded by $3C_AP,8P,16P,16P$, and
$N\kappa_1\le4P$. Hence
\[
 2D_1+N\kappa_1\le(6C_A+84)P\le D_p,\qquad
 4D_1<c_4:=c_*+2Nc_{\rm cl}+2n'c_{\rm cl}.           \tag{NP.8}
\]
The strict inequality uses $8(D_p+1)\le c_*$.
Thus $4D_0<c_4$ and $2D_1<c_4$ as required.

For $J_0=\lambda(a)b'/k$, $J=\lfloor J_0\rfloor$, the hypotheses give
$3\le J\le J_0\le b'/216$ and $c_{\rm cl}\le J_0$.
Let $u_a=\lfloor a_r+a_r/(kY)\rfloor$, $u_b=\lfloor b_u\rfloor$,
$l_b=\lceil b_l\rceil$. Then
\[
 \begin{gathered}
 J\le a,l_b,\qquad
 a+a/(kY)\le u_a+J+c_{\rm cl}\le a+a/(kY)+2J_0,\\
 u_b+J+c_{\rm cl}\le b_u+2J_0,\qquad
 (u_a+J+c_{\rm cl})+(u_b+J+c_{\rm cl})+1\le8a.
 \end{gathered}                                      \tag{NP.9}
\]
Each follows by substituting the floor/ceiling bounds; the report gives
the inequalities explicitly. Moreover $c_4\le2(c_*+Nc_{\rm cl}+n'c_{\rm cl})$
gives $\mathrm{CBoundI}(4)$, and
\[
 m_0=C_mkN_1Qc_4+12c_4\le b'/(kQY).                 \tag{NP.10}
\]
For (NP.10), multiply by $kQY$ and dominate the result by
$(2C_m+24)k^2N_1Q^2Y(c_4/2)$; the final size hypothesis applies.
Equations (NP.6)--(NP.10) supply exactly both
$\mathrm{FixedWeightNumericalFacts}$ and $\mathrm{MeanKappaNumericalFacts}$,
uniformly in all data with $C_H\ge T_M$.

\paragraph{NP7--NP8: the true final height.}
Every odd marked point $y$ belongs to exactly one open cyclic interval
between consecutive even marked points. Its segment start is
$y_{\rm prev}$ and its segment distance is $d_{\rm prev}$.
This is a bijection, including the interval crossing zero, so the
segment order sum equals
\[
 S=\sum_{y\ {\rm odd}}\min\{b_y,a_{y_{\rm prev}}/p^{d_{\rm prev}}\}.
\]
For $c\ge0,P\ge1$,
$\min(b+c,(a+c)/P)\le\min(b,a/P)+c$.
Thus shifting both degree families by $c_{\rm cl}$ changes $S$ to
$S'\le S+Nc_{\rm cl}$.
The actual integer height is $t_*=\lceil t_0+NJ_0\rceil$,
so it costs at most one more than $t_0+NJ_0$.
In comparing the local error terms at $a$ with the final ones at $a_r$,
their respective factors are $2,1,4,2$; the extra $NJ_0$ uses another
factor two on $N\lambda(a_r)b'/k$.
The allowance $CNkQ$ absorbs the rounding unit, since $NkQ\ge1$,
and $CNc_{\rm cl}$ absorbs $S'-S$. Therefore
\[
 C\ge4C_t+4\quad\Longrightarrow\quad
 t_*(a_w+c_{\rm cl},b_w+c_{\rm cl})+2n'c_{\rm cl}
       \le t_I(C_0,C;a_w,b_w).                      \tag{NP.11}
\]
No floor or additional shift has been suppressed.
Finally the size inequality also gives
$12Nkc_{\rm cl}/b'\le1/9$, $2c_{\rm cl}\le8a$,
and $b'/3\le a,b_l$. Together with (NP.8), these are precisely
$\mathrm{FactsLP}(K,C)$, supplied with thresholds
$(C_H,C)\ge(T_{\rm LP},4C_t+4)$.
All constant choices precede the variable data. Combining several such
packages only takes a maximum of their thresholds, and therefore does
not reverse any of the source quantifiers.

\paragraph{NP9: the explicit coarse loss constant.}
The actual loss is
\[
 \kappa_P=C_\kappa N\left(e^{-\lambda/8}
       +\frac{kQ(m_0+N_1k)}{b'}+\frac1Y+\frac{kE_0}a\right).
\]
Under NP2, set $\alpha=a/Y$ and $\beta=b'/Y$.
The standard size bounds give
$k^3Q^2c\le\beta$ and $k^3Q\le\beta$, while
$\lambda\le k$, $kb'\le a$, and $\lambda Yb'\le a$.
Since $\lambda\ge8\log(NY)$, the exponential term is at most $N_1/Y$.
Expanding $m_0=C_mkN_1Qc+12c$ bounds the second term by
$(C_m+13)N_1/Y$.
After multiplication by $k$, the five summands in $E_0$ are bounded
by $\alpha,\alpha,N_1\alpha,\alpha,N_1\alpha$, respectively.
Thus
\[
 \boxed{\quad\kappa_P\le C_s\,\frac{NN_1}{Y},\qquad
 C_s=C_\kappa(C_m+4C_{E1}+C_{E2}+15).\quad}          \tag{NP.12}
\]
The constant is chosen after $C_\kappa$ but before any variable data,
as required by the source's $\mathrm{FamilyFourthBound}$.


